To identify a concrete, data-backed entry point rather than an
a-priori guess, we read the action/observation traces of every
Assignment-2 $n{=}10$ ReAct episode on both domains.

\textbf{FEVER: premature confidence, invisible to the paper's trigger.}
Eight of ten ReAct episodes terminated after exactly one \texttt{Search}
followed immediately by \texttt{Finish}. The paper's trigger --- step
exhaustion without a \texttt{Finish} --- fired on \emph{none} of these,
because the model always did commit to an answer; it simply committed too
early. On the six episodes with gold label \texttt{NOT ENOUGH INFO}, the
model predicted \texttt{REFUTES} in all six: presented with a single,
often only tangentially relevant, retrieved passage, the model appears to
default to a confident negative rather than expressing uncertainty.

\textbf{HotpotQA: the paper's trigger works as intended.} By contrast,
every failing HotpotQA ReAct episode we inspected burned the full 7-step
budget without reaching \texttt{Finish}; one episode reissued
\texttt{Search[Hardley Flood]} verbatim after already having loaded that
page. Here the paper's trigger fires correctly.

This asymmetry is the entry point for this paper: the trigger is not
wrong, but it is \emph{incomplete}. It encodes only the failure mode most
salient in multi-hop QA (searching and re-searching without converging)
and has no mechanism for a claim-verification-style failure (converging
too fast on too little). Because both are genuine ReAct failure modes,
and the paper's own framing (Section 3.3) presents the trigger as a
general recipe rather than a HotpotQA-specific one, we treat this as a
limitation of the general hybrid mechanism rather than of the paper's
HotpotQA results specifically.
